\documentclass[11pt,a4paper]{article}

\usepackage[utf8]{inputenc}
\usepackage[T1]{fontenc}
\usepackage[margin=2.4cm]{geometry}
\usepackage{amsmath,amssymb}
\usepackage{booktabs}
\usepackage{array}
\usepackage{enumitem}
\usepackage{xcolor}
\usepackage{hyperref}
\hypersetup{colorlinks=true,linkcolor=black,citecolor=black,urlcolor=blue!50!black}

\setlength{\parskip}{5pt}
\setlength{\parindent}{0pt}

\title{DeltaFold --- Plan d'expérimentation (état actuel)\\[2pt]
\large Clustering structural non supervisé du virome par apprentissage topologique}
\author{Matthieu Nardi --- stage M2}
\date{\today}

\begin{document}
\maketitle

\begin{center}
\fbox{\parbox{0.92\linewidth}{\small
Document de travail. Il fixe les \emph{questions} auxquelles la prochaine campagne
doit répondre, puis le dispositif (données, objectif, architecture, métriques,
témoins, portes de décision) conçu pour y répondre. Il fait suite au post-mortem de
la première run, qui n'apprenait rien au-delà du signal des features.}}
\end{center}

\tableofcontents
\vspace{4pt}
\hrule

\section{But et cadrage}

\paragraph{But scientifique.} Reproduire, puis si possible étendre, l'étude de
Nomburg et al.\ (2024) sur le virome : recouvrer une organisation structurale des
protéines virales et, à terme, proposer de nouveaux analogues structuraux pour les
protéines « sombres » (sans homologue connu). L'outil est un encodeur structural
appris sur le complexe combinatoire de la protéine (hiérarchie
résidu $\to$ élément de structure secondaire $\to$ protéine).

\paragraph{Décisions actées.}
\begin{itemize}[leftmargin=1.4em,itemsep=2pt]
  \item \textbf{Non supervisé pur} : on retire le contrastif supervisé (positifs
        « même cluster Foldseek ») et le terme auxiliaire TM-score, pour éliminer
        toute fuite côté entraînement.
  \item \textbf{Architecture équivariante} (SE(3)), de type perceptron géométrique,
        consommant les vecteurs de déplacement C$\alpha$ et non plus seulement des
        scalaires invariants.
  \item \textbf{Objectif refondu} (voir §\ref{sec:dispositif}) : positifs par
        sous-structures, file de négatifs, tête de projection.
  \item \textbf{Base de prototypage corrigée} : représentative (avec singletons,
        distribution de tailles préservée), à distinguer du jeu complet.
\end{itemize}

\section{Les questions auxquelles le plan doit répondre}
\label{sec:questions}

\begin{description}[leftmargin=0pt,itemsep=4pt]
  \item[\textbf{Q1 --- Apport de la hiérarchie topologique.}] Les complexes
        combinatoires (hiérarchie résidu/SSE/protéine) encodent-ils quelque chose que
        les recherches structurales actuelles (Foldseek, embeddings globaux) ne
        capturent pas, en particulier l'\emph{homologie partielle / de domaine} ?
  \item[\textbf{Q2 --- Objectif et création des paires.}] Quelle stratégie de
        positifs et quel objectif permettent un \emph{apprentissage réel}, alors que
        la première run dégradait le signal des features au fil de l'entraînement ?
  \item[\textbf{Q3 --- Collapse.}] Comment éviter l'effondrement de la
        représentation observé en première run (sur les configurations à négatifs
        durs) ?
  \item[\textbf{Q4 --- Circularité de l'évaluation.}] Comment évaluer la qualité sans
        circularité vis-à-vis de Foldseek (qui est à la fois l'outil à battre et la
        « vérité terrain ») ?
  \item[\textbf{Q5 --- Géométrie : invariance vs équivariance.}] L'équivariance et la
        complétude géométrique apportent-elles, par rapport à un encodeur invariant
        scalaire ?
  \item[\textbf{Q6 --- Représentativité des données.}] La base de prototypage est-elle
        représentative, et les résultats transfèrent-ils au jeu complet ?
  \item[\textbf{Q7 --- Découverte (Nomburg).}] Peut-on retrouver les résultats de
        Nomburg et proposer de nouveaux analogues dans la fraction « sombre » ?
\end{description}

\section{Post-mortem de la première run (motivation)}

Sur la base de prototypage, un balayage d'ablations a montré que, pour presque toutes
les configurations, le meilleur TM-$\rho$ est atteint à l'epoch 1--2 puis se
\emph{dégrade}, pendant que la perte de validation baisse, et \emph{aucune}
configuration ne bat la baseline à poids aléatoires. Les runs « features/géométrie
seules » sont parmi les meilleurs. Diagnostic : l'entraînement contrastif ne fait pas
qu'échouer à apprendre, il \emph{détruit} le signal structural présent dans les
features d'entrée. Causes principales identifiées dans la revue de littérature :
\begin{enumerate}[leftmargin=1.5em,itemsep=2pt]
  \item \textbf{Positifs triviaux} : deux vues jitter/masquage de la protéine entière
        sont quasi identiques (perte par batch $\sim 10^{-5}$) ; le modèle apprend une
        invariance au bruit sans rapport avec la similarité de fold.
  \item \textbf{Négatifs durs $\to$ faux négatifs et collapse} : les quatre runs en
        collapse utilisent tous le \emph{hard-negative mining} ; tous les runs sans
        négatifs durs l'évitent et obtiennent les meilleurs scores.
  \item \textbf{Architecture en aval du goulot} : un encodeur plus expressif ne peut
        rien extraire d'un signal d'entraînement dégénéré.
\end{enumerate}

\section{Dispositif expérimental}
\label{sec:dispositif}

\subsection{Données}
Base de prototypage \emph{corrigée} ($\sim 3$--$5$k protéines) : conserver une part
de singletons, préserver la distribution des tailles de clusters (échantillonnage
proportionnel, pas « 2 par cluster »), retirer la sélection par mots-clés du jeu
d'entraînement, et adopter un découpage train/val « froid » (par cluster structural)
avec déduplication des quasi-doublons. Re-run sur le jeu complet exigé avant toute
conclusion biologique.

\subsection{Objectif refondu (trois composants)}
\begin{itemize}[leftmargin=1.4em,itemsep=3pt]
  \item \textbf{Positifs par sous-structures} (Hermosilla \& Ropinski) : pour chaque
        protéine, deux sous-structures connexes \emph{distinctes} (segment contigu par
        défaut ; boule spatiale en option). La taille est l'hyperparamètre critique
        (« sweet spot » : ni vues identiques, ni vues ne partageant plus le fold).
  \item \textbf{File de négatifs (MoCo)} : encodeur \emph{online} + encodeur
        \emph{momentum} (moyenne mobile, sans gradient) + file FIFO de clés négatives,
        ce qui découple le nombre de négatifs de la taille du batch. On \emph{retire}
        le hard-negative mining.
  \item \textbf{Tête de projection} : la perte opère sur la projection ; l'embedding
        d'\emph{évaluation} est pris \emph{avant} la tête (jetée), ce qui prévient le
        collapse dimensionnel.
\end{itemize}

\subsection{Architecture et ablations}
Encodeur équivariant sur le complexe combinatoire. Ablations, un facteur à la fois :
\begin{center}\small
\begin{tabular}{@{}p{3.2cm} p{5.0cm} p{6.0cm}@{}}
\toprule
Axe & Conditions & Question visée \\
\midrule
Hiérarchie & PCC complet vs $-$rang 2 vs $-$rang 3 (plat) vs GNN plat & Q1 \\
Géométrie & invariant scalaire vs GVP (norme) vs GCP (repères) & Q5 \\
Objectif & avec/sans tête ; MoCo vs sans négatifs (VICReg/Barlow) & Q2, Q3 \\
Positifs & taille de sous-structure ; contigu vs spatial & Q2 \\
Encodage distance & sinuso\"idal vs RBF & --- \\
\bottomrule
\end{tabular}
\end{center}

\section{Métriques}

\begin{itemize}[leftmargin=1.4em,itemsep=2pt]
  \item \textbf{Santé / collapse (porte bloquante)} : écart-type d'embedding, cosinus
        moyen hors-diagonale, rang effectif, alignement \& uniformité.
  \item \textbf{Accord structural sans clustering} : TM-recall@$k$ et corrélation
        distance--TM contre une vérité \textbf{TM-align} (et non Foldseek).
  \item \textbf{Clustering vs références indépendantes} : ARI, V-measure,
        homogénéité/complétude, Fowlkes--Mallows contre \textbf{CATH} (Foldseek en
        simple sanity check) ; indices de fragmentation/fusion.
  \item \textbf{Spécifique hiérarchie} : récupération d'homologie \emph{partielle} là
        où les méthodes globales échouent ; delta d'ablation des rangs.
  \item \textbf{Cohérence fonctionnelle (externe)} : pureté des clusters vis-à-vis
        EC/GO/Pfam.
  \item \textbf{Découverte} : nombre d'analogues « sombres » confiants (avec contrôle
        de FDR) ; contrôles positifs de Nomburg (RNA ligase T).
\end{itemize}

\section{Baselines et témoins}
Plancher : embedding aléatoire, modèle non entraîné. À comparer : Foldseek (global et
local), méthodes d'embedding (FoldExplorer, TM-Vec), GNN plat (GearNet, GVP),
décomposition en domaines $+$ Foldseek, U-Nets géométriques. Témoins : permutation des
clusters (null pour le supervisé), embedding mélangé.

\section{Correspondance question $\to$ test $\to$ critère}

\begin{center}\small
\begin{tabular}{@{}p{1.0cm} p{7.2cm} p{6.4cm}@{}}
\toprule
Q & Test / dispositif & Critère de décision \\
\midrule
Q1 & Ablation des rangs + test d'homologie partielle vs baselines forts & Le PCC complet bat $-$rang 2 / plat / (domaines$+$Foldseek) sur l'homologie partielle, sinon la hiérarchie n'apporte rien \\
Q2 & Positifs par sous-structures + MoCo + tête ; suivi loss/TM & Loss par batch non triviale \emph{et} TM-rho/recall montent avec l'entraînement (au lieu de dégrader) \\
Q3 & Tête de projection + file MoCo (repli VICReg/Barlow) & Rang effectif stable, cosinus moyen bas \\
Q4 & Évaluation vs CATH, TM-align continu, fonction (EC/GO/Pfam) & Accord avec des références \emph{indépendantes} de Foldseek, au-delà d'une distillation \\
Q5 & Échelle invariant $<$ GVP $<$ GCP & Le gain géométrique justifie le coût \\
Q6 & Base corrigée + re-run jeu complet & Les conclusions tiennent hors prototypage \\
Q7 & Encoder viral $+$ non-viral, retrieval multi-échelle, FDR & Reproduit les cas Nomburg ; analogues sombres sous seuil de FDR \\
\bottomrule
\end{tabular}
\end{center}

\section{Portes de décision (falsifiables)}
\begin{enumerate}[leftmargin=1.5em,itemsep=2pt]
  \item \textbf{Tâche non triviale} : si la perte par batch s'effondre ($\sim 10^{-5}$),
        réduire la taille des sous-structures.
  \item \textbf{Apprentissage réel} : si TM-rho/recall ne s'améliorent pas avec
        l'entraînement, l'objectif est encore mal posé.
  \item \textbf{Non-circularité} : si l'accord avec CATH/fonction n'excède pas une
        distillation de Foldseek, pas de contribution démontrée.
  \item \textbf{Hiérarchie} : si le PCC complet n'apporte rien sur l'homologie
        partielle, pivoter ou publier le résultat négatif.
\end{enumerate}

\section{Risques principaux}
Mémoire (deux encodeurs équivariants sur 16~Go) ; réglage de la taille des
sous-structures (nouveau point de défaillance) ; couverture CATH/fonction faible pour
le virome (d'où TM-align continu en complément) ; collapse persistant (repli sur un
objectif sans négatifs).

\end{document}
